\documentclass[10pt,a4paper]{amsart}
\usepackage[margin=19mm]{geometry}
\usepackage{amsmath,amssymb,mathtools}
\usepackage[scr=rsfs,cal=euler]{mathalfa}
\usepackage{xfrac}
\usepackage[unicode,hidelinks,bookmarks=false]{hyperref}
\newcommand{\ie}{i.e.\ }
\newcommand{\cf}{cf.\ }
\newcommand{\resp}{respectively\ }
\newcommand{\Subsection}{\subsection}
\providecommand{\nobreakdash}{\nobreak\hbox{-}}
\setlength{\emergencystretch}{2em}
\multlinegap0pt
\usepackage{bm}
\usepackage{bbm}
\usepackage{dsfont}
\usepackage{xspace}
\usepackage{cancel}
\usepackage{mathtools}
\usepackage{amsthm}
\makeatletter
\@ifundefined{thm}{\newtheorem{thm}{Theorem}}{}
\makeatother
\title[Lifting $I$-functions from the Grassmannians to their cotang]{Lifting $I$-functions from the Grassmannians to their cotangent bundles}
\author{Kamyar Amini}
\date{Source: arXiv:2511.06155v1 (CC BY 4.0)}
\begin{document}
\maketitle

\noindent\emph{Source and licence.} Lifting $I$-functions from the Grassmannians to their cotangent bundles. Version arXiv:2511.06155v1 (CC BY 4.0), distributed under the Creative Commons Attribution 4.0 International licence, as recorded in the arXiv licence metadata for this exact version and shown on the abstract page https://arxiv.org/abs/2511.06155v1. Full rights notice: licenses/arxiv-2511.06155-CC-BY-4.0.txt.\par\smallskip

\noindent\textit{Editorial scope.} Counted unit: the Thm whose source label is \texttt{main thm} in the article (source0.tex lines 1538--1548, source label \texttt{main thm}), delivered together with the article's own complete original proof of it (source0.tex lines 1549--1632, one \texttt{proof} environment of 84 lines). No mathematics is written, repaired or re-derived: statement and proof are the article's own text line for line. Required context retained uncounted: none. Every definition, display and label of those lines is retained unchanged. Omitted from this package and not redistributed: 1--1537, 1633--1918. Nothing inside a retained excerpt is omitted or altered: every retained body is delivered exactly as the article's lines are written, so no omission receipt of any kind accompanies this package. Citations: the retained excerpts carry no citation command at all, so no bibliography entry of the article is transcribed and no reference is redistributed. Reference adaptation: none was needed and none was made; every reference inside the delivered unit resolves inside it, so no unresolved reference or citation is produced. Printed slips kept literal: no printed word, symbol, hypothesis or number of the counted statement or of its proof was altered. Layout and class: the article's own class is \texttt{amsart} with packages including amsmath, amsthm, amsfonts, amssymb, amscd, framed, xy, nicefrac, color, graphicx, xcolor, tikz, tikz-cd, hyperref; the wrapper is the lane's pinned \texttt{amsart} editorial wrapper at 10pt on A4, which restates the theorem-like environments the retained text needs and adds the article's own macro definitions that the retained text needs (none), with extra wrapper packages (bm, bbm, dsfont, xspace, cancel, mathtools). The delivered numbering is the unit's own. Assets: no retained span contains a figure environment, a graphics-inclusion command, a PDF-page inclusion, a table or a TikZ picture; no image file is copied into this package, the package declares no asset dependency and no image file is redistributed with it. Licence: version arXiv:2511.06155v1 of this article is distributed under the Creative Commons Attribution 4.0 International licence, as recorded in the arXiv licence metadata for this exact version and shown on the abstract page https://arxiv.org/abs/2511.06155v1. A full rights notice is delivered at licenses/arxiv-2511.06155-CC-BY-4.0.txt. Characters: every retained excerpt is pure ASCII, so the delivered engine, XeLaTeX with its default Latin Modern text font, reports no missing character.
\begin{thm} \label{main thm}
     After the change of variables $y=-q^{-1}\hbar$, we have: 
    \[
    V_p(t;q,Q,\hbar) = \mathcal{B}_{-q^{-1}\hbar}(I(t;q,Q)_{\arrowvert_{p}})
    \]

    where the coefficients of $Q^d$ in both functions are equal for every $d \geq 0$, namely,
    \[
    V_p^d(t;q,\hbar) = \mathcal{B}_{-q^{-1}\hbar}(I_d(t;q)_{\arrowvert_{p}}).
    \]
\end{thm}

\begin{proof}
     In the expression for $V_p^d$, we have terms of the form $\prod_{i=1}^{r}\prod_{j=1}^{n}\{\frac{t_j}{t_i}\}_{d_i}$. We divide this term into two parts:

   \begin{itemize}
       \item The case $i=j$:
       
       In this case, we have $r$ terms of $\{1\}_{d_i}$ for $1 \leq i \leq r$. Since $\{1\}_{d_i} = \frac{(1-\hbar)\dots(1-\hbar q^{d_i-1})}{(1-q)\dots(1-q^{d_i})}$, we obtain:
       \[
       \frac{\prod_{i=1}^r \prod_{s=0}^{d_i-1}(1-\hbar q^s)}{\prod_{i=1}^r \prod_{s=1}^{d_i}(1-q^s)}.
       \]
       Now consider $\mathcal{B}_{y}(I(t;q,Q)_{\arrowvert_{p}})$, which contains the terms: 
       $\frac{\prod_{i=1}^r\prod_{s=1}^{d_i}(1+yq^{s})}{\prod_{i=1 }^r\prod_{s=1}^{d_i}(1-q^{s})}$.
       
       The corresponding term in $\mathcal{B}_{q^{-1}y}(I(t;q,Q)_{\arrowvert_{p}})$ will be:
       \[
       \frac{\prod_{i=1}^r\prod_{s=1}^{d_i}(1+yq^{s-1})}{\prod_{i=1 }^r\prod_{s=1}^{d_i}(1-q^{s})} = \frac{\prod_{i=1}^r\prod_{s=0}^{d_i-1}(1+yq^{s})}{\prod_{i=1 }^r\prod_{s=1}^{d_i}(1-q^{s})}.
       \]
       
       Hence, after the change of variable $y=-\hbar$, the terms are equal in $V_p^d$ and $\mathcal{B}_{-q^{-1}\hbar}(I(t;q,Q)_{\arrowvert_{p}})$.

       \item The case $1 \leq i \leq r$ and $r+1 \leq j \leq n$.
       
       These terms in $V_p^d$ are of the form: $\prod_{i=1}^{r}\prod_{j=r+1}^{n}\{\frac{t_j}{t_i}\}_{d_i}$, and we have:
       \[
       \{\frac{t_j}{t_i}\}_{d_i} = \frac{(1-\hbar \frac{t_i}{t_j})\dots(1- \hbar q^{d_i-1}\frac{t_i}{t_j})}{(1-q \frac{t_i}{t_j})\dots(1-q^{d_i}\frac{t_i}{t_j})}.
       \]
       Thus, the overall expression in $V_p^d$ becomes:
       \[
\prod_{i=1}^r \prod_{j=r+1}^{n} \frac{(1-\hbar \frac{t_i}{t_j})\dots(1- \hbar q^{d_i-1}\frac{t_i}{t_j})}{(1-q \frac{t_i}{t_j})\dots(1-q^{d_i}\frac{t_i}{t_j})}   
       \]

       Now, we examine $\mathcal{B}_{-q^{-1}\hbar}(I(t;q,Q)_{\arrowvert_{p}})$. In $\mathcal{B}_{y}(I(t;q,Q)_{\arrowvert_{p}})$, we have the terms: 
       \[
       \frac{\prod_{i=1}^r \prod_{j=r+1}^n\prod_{s=1}^{d_i}(1+yq^{s}\frac{t_i}{t_j})}{\prod_{i=1}^r \prod_{j=r+1}^n\prod_{s=1}^{d_i}(1-q^{s}\frac{t_i}{t_j})} 
       \]
       Hence, the corresponding term in $\mathcal{B}_{q^{-1}y}(I(t;q,Q)_{\arrowvert_{p}})$ will be:

       \[
       \frac{\prod_{i=1}^r \prod_{j=r+1}^n\prod_{s=1}^{d_i}(1+yq^{s-1}\frac{t_i}{t_j})}{\prod_{i=1}^r \prod_{j=r+1}^n\prod_{s=1}^{d_i}(1-q^{s}\frac{t_i}{t_j})} = \frac{\prod_{i=1}^r \prod_{j=r+1}^n\prod_{s=0}^{d_i-1}(1+yq^{s}\frac{t_i}{t_j})}{\prod_{i=1}^r \prod_{j=r+1}^n\prod_{s=1}^{d_i}(1-q^{s}\frac{t_i}{t_j})}
       \]
       Thus, after the change of variables $y=-\hbar$, the terms in $V_p^d$ and $\mathcal{B}_{-q^{-1}\hbar}(I(t;q,Q)_{\arrowvert_{p}})$ are also equal.

       \item Finally, in $V_p^d$, we are left with the terms: $\prod_{i,j=1}^r\{\frac{t_j}{t_i}\}_{d_i}$ where $1 \leq i \ne j \leq r$, and the terms $\prod_{i,j=1}^r\{\frac{t_j}{t_i}\}_{d_i-d_j}^{-1}$. 
       
       In $\prod_{i,j=1}^r\{\frac{t_j}{t_i}\}_{d_i-d_j}^{-1}$, when $i=j$, we have terms of the form $\{1\}_0$, which are equal to 1. So we can focus on the case when $i \ne j$. We then obtain the following:
       \[
\prod_{i,j=1}^r\{\frac{t_j}{t_i}\}_{d_i}\prod_{1 \le i,j \le r, i \ne j}\{\frac{t_j}{t_i}\}_{d_i-d_j}^{-1} = \prod_{1 \le i,j \le r, i \ne j}\{\frac{t_j}{t_i}\}_{d_i-d_j}^{-1}\{\frac{t_j}{t_i}\}_{d_i}.
       \]
       Now we calculate $\{\frac{t_j}{t_i}\}_{d_i-d_j}^{-1}.\{\frac{t_j}{t_i}\}_{d_i}$:
       
        1. If $d_i-d_j > 0:$

        \begin{align*}
           \{\frac{t_j}{t_i}\}_{d_i-d_j}^{-1}\{\frac{t_j}{t_i}\}_{d_i} &= \frac{\prod_{s=1}^{d_i-d_j}(1-q^s\frac{t_i}{t_j})}{\prod_{s=0}^{d_i-d_j-1}(1-\hbar q^s \frac{t_i}{t_j})} \cdot\frac{\prod_{s=0}^{d_i-1}(1-\hbar q^s\frac{t_i}{t_j})}{\prod_{s=1}^{d_i}(1-q^s \frac{t_i}{t_j})}  \\&= \frac{\prod_{s=d_i-d_j}^{d_i-1}(1-\hbar q^s \frac{t_i}{t_j})}{\prod_{s=d_i-d_j+1}^{d_i}(1-q^s \frac{t_i}{t_j})}
        \end{align*}
        2. If $d_i-d_j < 0:$

        \begin{align*}
            \{\frac{t_j}{t_i}\}_{d_i-d_j}^{-1}\{\frac{t_j}{t_i}\}_{d_i} &= \frac{\prod_{s=d_i-d_j}^{-1}(1-\hbar q^s  \frac{t_i}{t_j})}{\prod_{s=d_i-d_j+1}^0(1-q^s\frac{t_i}{t_j})} \cdot \frac{\prod_{s=0}^{d_i-1}(1- \hbar q^s \frac{t_i}{t_j})}{\prod_{s=1}^{d_i}(1-q^s\frac{t_i}{t_j})}  \\&= \frac{\prod_{s=d_i-d_j}^{d_i-1}(1-\hbar q^s  \frac{t_i}{t_j})}{\prod_{s=d_i-d_j+1}^{d_i}(1-q^s\frac{t_i}{t_j})}.
        \end{align*} 
       In both cases, the expressions are the same, so the total term is:
        \begin{align*}
            \prod_{1 \leq i,j \leq r}^{i \ne j}\frac{\prod_{s=d_i-d_j}^{d_i-1}(1-\hbar q^s \frac{t_i}{t_j})}{\prod_{s=d_i-d_j+1}^{d_i}(1-q^s\frac{t_i}{t_j})}
        \end{align*}
        
        We now turn our attention to $\mathcal{B}_{y}(I(t;q,Q)_{\arrowvert_{p}})$, where we are left with the following terms:
        \begin{align*}
        \frac{\prod_{1 \leq i,j \leq r}^{i \ne j}\prod_{s=-d_i}^{-d_i+d_j-1}(1+yq^{-s}\frac{t_i}{t_j})}{\prod_{1 \leq i,j \leq r}^{i \ne j}\prod_{s=-d_i}^{-d_i+d_j-1}(1-q^{-s}\frac{t_i}{t_j})}.
        \end{align*}
        This can be rewritten as:
        \begin{align*}
        \frac{\prod_{1 \leq i,j \leq r}^{i \ne j}\prod_{s=d_i-d_j+1}^{d_i}(1+yq^{s}\frac{t_i}{t_j})}{\prod_{1 \leq i,j \leq r}^{i \ne j}\prod_{s=d_i-d_j+1}^{d_i}(1-q^{s}\frac{t_i}{t_j})}.
        \end{align*}
        Therefore, the terms in $\mathcal{B}_{q^{-1}y}(I(t;q,Q)_{\arrowvert_{p}})$ are given by:
        \begin{align*}
        \frac{\prod_{1 \leq i,j \leq r}^{i \ne j}\prod_{s=d_i-d_j+1}^{d_i}(1+yq^{s-1}\frac{t_i}{t_j})}{\prod_{1 \leq i,j \leq r}^{i \ne j}\prod_{s=d_i-d_j+1}^{d_i}(1-q^{s-1}\frac{t_i}{t_j})},
        \end{align*}
            which are equal to terms:
            \begin{align*}
                \frac{\prod_{1 \leq i,j \leq r}^{i \ne j}\prod_{s=d_i-d_j}^{d_i-1}(1+yq^{s}\frac{t_i}{t_j})}{\prod_{1 \leq i,j \leq r}^{i \ne j}\prod_{s=d_i-d_j+1}^{d_i}(1-q^{s}\frac{t_i}{t_j})}.
            \end{align*}
            Thus, after the change of variables $y=-\hbar$, the terms in $V_p^d$ and $\mathcal{B}_{-q^{-1}\hbar}(I(t;q,Q)_{\arrowvert_{p}})$ are also equal.
    \end{itemize}
    \end{proof}

\end{document}
